考虑实值、足够光滑且在空间无穷远迅速衰减的无量纲 Korteweg–de Vries 方程
\[
u_t+6uu_x+u_{xxx}=0,\qquad x\in\mathbb R.
\]
所有积分都在整条实线上进行，假定边界项消失。
(1) 令 \(u(x,t)=U(\xi)\)，\(\xi=x-vt-x_0\)，并要求 \(U,U',U''\to0\) 于 \(|\xi|\to\infty\)。从行波常微分方程及第一积分出发，严格求出正幅值局域孤立波、波速和波峰振幅的关系。
(2) 从偏微分方程证明质量 \(M=\int u\,dx\) 与二次积分 \(I=\int u^2\,dx\) 守恒；对 (1) 的孤立波精确求这两个积分，并分析其对波速的依赖。
(3) 验证哈密顿泛函 \(H=\int[u^3-\tfrac12(u_x)^2]dx\) 守恒，并计算孤立波的 \(H\)。解释其中梯度能项的作用，核验各积分的尺度指数。
(4) 证明常数背景的 Galilei 变换及 KdV 尺度变换；线性化求色散关系与群速度，说明为何正向孤立波可以与线性波具有相反传播方向。
